\begin{problem}
    En una empresa hay \num{4} trabajadores para cosechar uvas y estos logran llenar \num{2000} cajas en \num{20} días de trabajo, todos trabajando todos los días, al mismo ritmo y la misma cantidad de horas.

    En la empresa se realiza el siguiente procedimiento para calcular la cantidad de trabajadores necesarios para llenar \num{32000} cajas en \num{10} días, bajo las mismas condiciones, cometiéndose un error.

    \begin{itemize}
        \item Paso 1: se calcula la cantidad de cajas que llenan diariamente los \num{4} trabajadores iniciales en conjunto como $\dfrac{2000}{20}$, obteniéndose que llenan \num{100} cajas diarias.
        \item Paso 2: se calcula la cantidad de cajas que cada trabajador llena diariamente como $100 \cdot 4$, obteniéndose \num{400} cajas diarias por trabajador.
        \item Paso 3: se calcula la cantidad de cajas que se deben llenar diariamente para la nueva producción como $\dfrac{32000}{10}$, obteniéndose que deben llenar \num{3200} cajas diarias.
        \item Paso 4: se divide la cantidad de cajas diarias que deben llenar por la cantidad que cada trabajador llena diariamente como $\dfrac{3200}{400}$, obteniéndose que necesitan \num{8} trabajadores.
    \end{itemize}

    ¿En cuál de los pasos se cometió el error?

    \begin{enumerate}[label=\Alph*.]
        \item En el Paso 1
        \item En el Paso 2
        \item En el Paso 3
        \item En el Paso 4
    \end{enumerate}
\end{problem}
